% !TeX root = ../main.tex
%------------------------------------------------------------------------------
%  Chapter 21 -- Mechanised Tunnelling
%  Sources: Maidl et al., Mechanised Shield Tunnelling and Hardrock TBMs;
%  Rostami & Ozdemir (CSM model); Bruland (NTNU model); Barton (Q-TBM);
%  DAUB recommendations for TBM selection.
%------------------------------------------------------------------------------
\chapter{Mechanised Tunnelling}
\label{ch:tbm}

Machine selection for mechanised tunnelling is committed at procurement, on a
ground model that is necessarily provisional, and cannot be revisited once the
drive has begun. A substantial proportion of the serious difficulties recorded
in mechanised tunnelling can be traced to a mismatch between the machine
specified and the ground subsequently encountered. This chapter therefore
treats selection as a geotechnical decision before a plant decision, and then
examines the interactions that govern performance: cutting tools against rock
abrasivity, thrust and torque against ground reaction, face pressure against
stability, and annulus grouting against ground movement.

This chapter therefore treats machine selection as a geotechnical decision
before a plant decision, and then works through the interactions that govern
performance: cutting tools against rock abrasivity, thrust and torque against
ground reaction, face pressure against stability, and grout against ground
movement.

\section{Machine selection and configuration}
\label{sec:machine-selection}

\subsection{Machine families}
\label{sec:machine-families}

The families divide by how, and whether, the machine supports the ground.

\emph{Open gripper} machines work in self-supporting rock. They brace against
the tunnel walls to react thrust, leaving the crown accessible so that rock
bolts, mesh and shotcrete can be installed close behind the head. They give
the geotechnical engineer what no shielded machine does --- a visible tunnel
wall --- and they cannot cope with ground that will not hold a grip pad.

\emph{Single shield} machines are enclosed, react thrust against the completed
segmental ring, and support the perimeter as they advance. \emph{Double
shield} machines add a gripper section so that boring and ring building can
proceed simultaneously, giving the highest advance rates in good rock while
retaining the ability to work in poor ground on the single-shield principle.
The cost is length and mechanical complexity, and a long machine negotiates
tight curves and squeezing ground poorly.

\emph{Closed-face} machines --- earth pressure balance and slurry ---
pressurise the excavation chamber and support the face as well as the
perimeter. They are required wherever the face will not stand, which in
practice means most soft ground and most work below the water table in
permeable material.

Selection follows the ground and water rather than preference: rock strength
and abrasivity, discontinuity spacing, the grading and permeability of any
soil, the water head, the presence of mixed faces, and the settlement limits
imposed from the surface. Alignment length and curvature then determine
whether the economics favour the higher-output configuration.

\subsection{Closed-face machines}
\label{sec:closed-face}

An \emph{earth pressure balance} machine supports the face with the excavated
soil itself, held in the chamber as a plastic paste and discharged at a
controlled rate through a screw conveyor. The screw is the pressure control:
material passing along it loses pressure gradually, so a plug of conditioned
soil maintains the differential between chamber and atmosphere. The
requirement is a paste of the right consistency --- soft enough to flow,
stiff and impermeable enough to form a plug --- and natural soils rarely
supply it. Conditioning with foam, polymer or bentonite is therefore normal
rather than exceptional, and the applicability envelope is defined by grading:
clays and silts condition readily, uniform sands and gravels with difficulty
or not at all.

A \emph{slurry} machine supports the face with pressurised bentonite
suspension, with an air bubble in a rear compartment regulating the pressure
precisely. Excavated material is removed hydraulically and separated at a
surface plant, and it is the separation plant, not the cutterhead, that
usually limits advance in fine soils: separating silt and clay from bentonite
is slow, and a plant sized for the average will stop the machine at the peak.

Support transfer differs fundamentally between them, and this is where slurry
machines fail. Pressure reaches the soil skeleton through a filter cake formed
as the slurry penetrates a short distance and its solids bridge the pore
openings. In coarse or gap-graded ground the slurry penetrates instead of
sealing, and the support pressure is dissipated into the ground rather than
applied to the face. Coarser bentonite, additives and increased density extend
the range; beyond it, the machine cannot support the face at all.

\emph{Multi-mode} machines can be converted between modes, and alignments
crossing several ground types increasingly justify them. Conversion is not
instantaneous --- it is a planned operation of days --- so the value lies in
crossing a small number of long, well-identified domains rather than in
responding to variability at the face.

\subsection{Cutterhead and cutting tools}
\label{sec:cutterhead-tools}

The cutterhead carries the tools, delivers thrust and torque to the face, and
admits muck to the chamber. Its opening ratio expresses the trade at the heart
of its design: a large opening ratio admits muck freely and gives poor face
support, while a closed head supports well and risks clogging in cohesive
ground. Rock machines therefore run open heads and closed-face machines run
substantially closed ones.

Hard rock is cut by disc cutters, which indent the rock and propagate cracks
between adjacent tracks so that chips spall out between them. Efficiency
depends on the spacing-to-penetration ratio: too close and the machine grinds
rather than chips, too far and the cracks fail to link. Soft ground is cut by
drag picks and scrapers, which shear rather than indent.

Wear governs cost and programme more than cutting does. Abrasivity indices ---
the Cerchar index for rock, the LCPC index for soils --- correlate with wear
rate and should be measured during investigation, because they determine how
often the machine must be stopped and the face entered. In hard abrasive rock
cutter changes can consume more time than boring, and in closed-face machines
under pressure a change becomes a hyperbaric intervention with its own safety
regime and its own days of downtime.

\section{Performance, pressure and control}
\label{sec:performance-pressure}

\subsection{Performance and installed power}
\label{sec:performance-power}

Two prediction models dominate. The Colorado School of Mines approach~\citep{rostami1993csm} is
mechanistic: it estimates the forces on an individual disc cutter from rock
properties and cutter geometry, sums them over the head, and returns thrust,
torque and penetration per revolution. The Norwegian University of Science and
Technology approach is empirical, built from an extensive case base and
correlating penetration with rock drillability and fracturing indices. Barton's
$Q_{TBM}$ extends rock mass classification to machine performance.

All three predict \emph{penetration rate} --- advance while boring. None
predicts \emph{advance rate}, which is what the programme needs, because
advance rate is penetration rate multiplied by utilisation, and utilisation is
the fraction of time the machine is actually boring. Utilisation on a
well-run drive in good ground may reach 40 to 50 per cent; in difficult ground
it can fall below 20. The remainder is ring building, cutter changes,
maintenance, logistics and stoppages.

This is why predicted advance rates are so often optimistic: the penetration
prediction is reasonable and the utilisation assumption is not. Realistic
programming requires utilisation drawn from comparable projects rather than
from the machine supplier.

Thrust must overcome cutting force at the face, friction along the shield, and
drag from the back-up. Torque must overcome cutterhead cutting resistance and,
in closed-face machines, the resistance of the material in the chamber.
Under-specification of either is unrecoverable: a machine without the torque to
turn its head in sticky clay, or without the thrust to overcome a squeezing
shield, has to be freed by intervention from outside.

Shield jamming in squeezing ground is the characteristic hard-rock failure
mode. Convergence closes onto the shield, friction rises beyond available
thrust, and the machine stops. Prediction requires the squeezing assessment of
Chapter~\ref{ch:high-stress}; mitigation is a shorter shield, a larger overcut,
higher installed thrust and the ability to overbore.

\subsection{Face pressure and annulus grouting}
\label{sec:face-pressure-grouting}

The operating pressure window follows directly from the analysis of
Chapter~\ref{ch:ground-behaviour}: a lower bound from face collapse plus water
pressure, an upper bound from blow-out and heave. Converting that window into
machine operating parameters requires care about where in the profile the
target is defined, since pressure varies over the face height with the density
of the supporting medium.

Control in real time is the operator's task and the designer's specification.
Chamber pressure responds to advance rate and discharge rate together, and
maintaining it through a ring-build stop --- when the machine is not advancing
but the ground is still there --- is a common point of loss.

Annulus grouting closes the gap between the excavated bore and the segment
ring, and it is the largest controllable component of volume loss. Simultaneous
grouting, injected through the shield tail as the machine advances, fills the
gap before the ground can converge into it; secondary grouting through the
segments is a repair rather than a method. Grout pressure must exceed the
ground and water pressure enough to fill the annulus and stay below the
fracture pressure, and the volume injected must be reconciled against the
theoretical annular volume ring by ring. A persistent shortfall between
theoretical and actual volume is the earliest reliable indicator of ground
loss.

\subsection{Guidance and interventions}
\label{sec:guidance-interventions}

Guidance systems track the shield continuously against the design alignment
and report deviation in position and attitude. Steering corrections are made
by differential thrust across the jack groups, and on articulated machines by
the articulation joint. Because correction is gradual, the build must
anticipate: universal ring taper is used to steer the ring as well as to
negotiate curves, and the build sequence is planned several rings ahead.

Roll, alignment drift and the accumulation of build tolerance all feed into
the clearance envelope, and a drive that consumes its tolerance early has no
margin left for the difficult ground it has not yet reached.

Interventions --- entering the excavation chamber for cutter changes,
inspection or repairs --- are the principal safety risk in closed-face
tunnelling. At low pressure, compressed air working is possible with medical
supervision and decompression control. Above the limits of compressed air,
saturation techniques are used, with their own substantial regime. The
alternative is to intervene under atmospheric conditions, which requires a
stable face --- achievable by locating interventions where the ground is
competent, or by treating the ground from the surface in advance.

Planning interventions is therefore an alignment-stage activity. Identifying
in advance where the ground allows a safe atmospheric intervention, and
scheduling cutter changes to those locations, is worth more than any amount of
capability to intervene under pressure.

\section{Machine data and forward exploration}
\label{sec:machine-data}

A tunnel boring machine records thrust, torque, penetration rate, cutterhead
rotation, chamber pressure, screw or slurry flow, grout volumes and pressures,
and muck weight, continuously and ring by ring. That record is the densest
ground data most projects ever collect, and it is generated free.

Interpreted as geotechnical information it supports three things: inference of
ground conditions behind the face, detection of departure from expected
behaviour in real time, and calibration of performance prediction for the
remainder of the drive. Specific energy --- the energy consumed per unit
volume excavated --- is the most useful derived quantity, since it responds to
rock strength and fracturing and can be tracked continuously against the ground
model of Chapter~\ref{ch:ground-modelling}.

The obstacles are practical. Data arrive at high rate in supplier-specific
formats, require cleaning and ring-based aggregation to be usable, and are
frequently retained by the contractor rather than shared with the designer.
Specifying access to machine data, in a defined format and at a defined
frequency, is a contract matter that should be settled before procurement.

Forward exploration is the other information source, and it is the only one
that reports ground the machine has not yet reached. Probe drilling ahead of a
shielded machine is slow and requires ports through the head, so it costs
directly in advance rate. Seismic and electromagnetic look-ahead methods are
faster and less definite. The decision to probe is a policy set from the
geological model --- probing through reaches identified as at risk of water,
fault or cavity --- rather than a response to something already observed,
because by the time a shielded machine observes a problem it is in it.

\section{Risk and logistics}
\label{sec:tbm-risk-logistics}

The characteristic failure modes are few and well documented. Shield jamming
in squeezing ground. Cutterhead clogging in sticky clay, where cohesive
material packs the openings and torque rises until the head stalls. Boulders
in soft ground, which damage tools and, in a closed-face machine, cannot be
seen before they are struck. Mixed faces, which produce uneven wear, steering
difficulty and over-excavation of the softer horizon
(Chapter~\ref{ch:ground-behaviour}). Karst cavities, which cause sudden loss
of face pressure and grout. Excessive wear in abrasive ground, which converts
a boring operation into a sequence of interventions.

Each has a ground signature that investigation can find and a mitigation that
must be specified before procurement: overcut and thrust for squeezing,
conditioning and head design for clogging, tool selection and rock crushers
for boulders, probing and grouting for karst.

Logistics constrain the drive as firmly as the geology. Launch requires a
shaft or portal, a cradle, a reaction frame and a seal against the ground; the
launch seal is a routine source of water and ground loss. Muck removal, segment
delivery, ventilation and power all run through a single tunnel behind the
machine, and their capacity sets a ceiling on advance that no improvement in
penetration rate can lift. The surface site --- for segment storage, slurry
separation, spoil handling and reception --- is often the binding constraint in
urban work, and it is fixed by land availability rather than by engineering.

Reception is planned at the same time as launch. A machine arriving at a shaft
must break through a face that is stable, be recovered or abandoned in place,
and leave a structure that can be completed. Machines are occasionally buried
where recovery is uneconomic, and that decision, if it is a possibility, is
better taken deliberately at design stage than under pressure at the end of a
drive.
